\begin{center}
Problema D

Agência de Espionagem

{\small Nome do arquivo: D.c, D.cpp, D.java, ou D.py}

{\large Timelimit: $1$}

{\tiny Por Athos José}


\end{center}

A Agência Nacional de Espionagem mantém em seus cofres uma coleção de símbolos secretos, representados por uma sequência de letras. Esses símbolos podem ser reorganizados em qualquer ordem, de forma a formar novas mensagens.

Agentes inimigos, porém, estão tentando criar mensagens escondidas dentro dessas reorganizações. Uma mensagem inimiga é considerada possível se puder aparecer como uma sequência contínua de caracteres de alguma permutação dos símbolos guardados no cofre.

Sua missão é, para cada mensagem inimiga recebida, determinar se ela pode ser formada a partir dos símbolos da Agência.

\vspace{0.3cm}
\textbf{Entrada}

A primeira linha contém um inteiro $N$ ($1 \leq N \leq 10^5$), o número de mensagens inimigas a serem analisadas, seguido de uma string $S$ ($1 \leq |S| \leq 10^5$) representando os símbolos secretos do cofre.

Cada uma das próximas $N$ linhas contém uma string $Q_i$ ($1 \leq |Q_i| \leq 10^5$), representando uma mensagem inimiga a ser verificada.

A soma dos comprimentos de todas as mensagens $Q_i$ não excede $10^9$.

Todas as strings são compostas apenas de letras minúsculas do alfabeto inglês (a-z).

\vspace{0.3cm}
\textbf{Saída}

Para cada mensagem $Q_i$ imprima uma linha contendo:
\begin{itemize}
    \item ``SIM'', se a mensagem pode ser formada como substring de alguma permutação de $S$.
    \item ``NAO'',caso contrário.
\end{itemize}

\begin{table}[h]
    \centering
    \begin{tabular}{|p{7cm}|p{7cm}|}
       \hline
        \begin{center}
            \vspace{-0.3cm}
            \textbf{Exemplos de Entrada}
            \vspace{-0.3cm}
        \end{center} 
         &  
        \begin{center}
            \vspace{-0.3cm}
            \textbf{Exemplos de Saída}
            \vspace{-0.3cm}
        \end{center} \\ \hline
         \texttt{3 abac} & \texttt{SIM} \\  
         \texttt{caa} & \texttt{NAO} \\
         \texttt{aaa} & \texttt{SIM} \\  
         \texttt{bac} & \\
         \hline
    \end{tabular}
    \caption{Exemplos de entradas e saídas}
    \label{tab:inputB}
\end{table}

\newpage

\begin{center}
\noindent
\lstinputlisting{abobora_25/D/D5.cpp}
\end{center}